% !TeX root = ./rpz.tex
\Introduction

Актуальность темы выпускной квалификационной работы обусловлена ростом числа повседневных и профессиональных задач, привязанных одновременно к географическому адресу и временному окну. Ручное планирование подобных маршрутов трудоемко и подвержено ошибкам, что формирует устойчивую потребность в инструментах автоматизированного планирования. Аналогичные тенденции наблюдаются и в коммерческой логистике последней мили~\cite{shkodinskiy2025lastmile}, где увеличение числа адресов доставки и сужение временных окон обслуживания делают задачу построения оптимального маршрута экономически значимой. Применение информационных технологий в управлении временем рассматривается как фактор повышения эффективности планирования личной и профессиональной деятельности~\cite{bronnikova2010tm}.

Существующие цифровые решения представлены либо приложениями для управления задачами, ориентированными на учет перечня дел без учета пространственного фактора, либо навигационными сервисами, формирующими маршрут без учета длительности выполнения задач. В результате пользователи вынуждены комбинировать несколько приложений, вручную сопоставляя адреса, интервалы времени и порядок посещения точек. Подобный подход характеризуется высокой трудоемкостью и ростом вероятности ошибок при увеличении количества задач~\cite{reshenie_zadach_planirovaniya2020}. Специализированные логистические платформы, напротив, ориентированы на крупный коммерческий сегмент и недоступны частному пользователю без коммерческой подписки.

Объект исследования~\textendash{} процесс планирования последовательности выполнения задач в городской среде с учетом пространственных и временных параметров.

Предмет исследования~\textendash{} алгоритмы маршрутизации и методы оптимизации порядка выполнения задач с временными окнами в условиях, характерных для частного пользователя и малого бизнеса.

Цель работы~\textendash{} разработка веб-приложения для автоматизированного формирования оптимальной последовательности выполнения задач с учетом географического положения, временных окон и длительности выполнения.

Для достижения поставленной цели сформулированы следующие задачи:

\begin{compactlist}
      \item проанализировать предметную область и существующие программные решения, выявить ограничения существующих инструментов;
      \item исследовать алгоритмические подходы к маршрутизации с временными окнами, обосновать выбор алгоритма;
      \item разработать архитектуру программной системы и модель данных;
      \item реализовать серверную часть и алгоритм оптимизации маршрута с поддержкой дополнительных ограничений;
      \item разработать пользовательский интерфейс для планирования задач, визуализации маршрута на карте, импорта и экспорта задач;
      \item реализовать интеграцию с внешними картографическими сервисами;
      \item провести модульное и нагрузочное тестирование разработанного решения;
      \item оценить производительность и масштабируемость системы.
\end{compactlist}

Научная и практическая новизна работы заключается в том, что предложенная реализация объединяет в одном инструменте управление задачами, автоматическую оптимизацию маршрута с учетом временных окон и попарных ограничений порядка выполнения, а также визуализацию маршрута. Сочетание перечисленных функций отсутствует у рассмотренных аналогов в сегменте бесплатных решений для индивидуального пользователя.

Практическая значимость работы заключается в том, что разработанное решение представляет собой единый инструмент планирования маршрута с временными окнами, доступный частным пользователям, специалистам на выезде, курьерам малого бизнеса и индивидуальным предпринимателям без коммерческой подписки. Веб-форма распространения и контейнеризация на основе Docker Compose обеспечивают воспроизводимое развертывание на типовом серверном оборудовании.

Методы исследования включают анализ предметной области и сравнительный анализ существующих решений, обзор и сравнение алгоритмических подходов к задаче маршрутизации с временными окнами (полный перебор, динамическое программирование, жадные эвристики, метаэвристики), методы объектно-ориентированного проектирования (принципы чистой архитектуры~\cite{martin2017clean}, предметно-ориентированное проектирование~\cite{evans2011ddd}, паттерны проектирования~\cite{gamma2009patterns}), модульное и нагрузочное тестирование~\cite{kulikov2017testing,myers2018testing}, экспериментальную оценку производительности.

Структура работы. Работа состоит из введения, трех глав, заключения, списка использованных источников и трех приложений. В первой главе выполнены анализ предметной области и обоснование выбора технологий. Во второй главе представлены проектные решения и реализация. В третьей главе рассмотрены вопросы внедрения, отказоустойчивости и тестирования. В приложениях приведены листинги SQL-миграций и конфигурации, листинги ключевых фрагментов программного кода, а также техническое задание на разработку.
